\begin{derivation}[从\eqnrefs{eq:sm-hgg-amplitude-dp94}到\eqnrefs{eq:sm-hgg-effective-operator-dp94}]
先对单个夸克味 $q$ 计算完整三角圈，而不预先代入重顶极限。Higgs--夸克约化顶角为 $-\ii m_qc^2/v_E$，夸克--胶子顶角为 $-\ii g_s\gamma^\mu T^a$。取两个出射胶子动量为 $k_1,k_2$，满足
\begin{equation*}
k_1^2=k_2^2=0,
\qquad
2k_1\!\cdot k_2=m_h^2c^2.
\end{equation*}
两种外胶子排序都必须保留。提出共同外腿归一化后，第一种排序的圈张量可写成
\begin{align*}
\mathcal M_{q,1}^{ab,\mu\nu}
={}&-g_s^2\frac{m_qc^2}{v_E}
\Tr(T^aT^b)
\int_\ell^{(d)}
\frac{\Tr\!\left[
(\slashed\ell+m_qc)\gamma^\mu
(\slashed\ell+\slashed k_1+m_qc)\gamma^\nu
(\slashed\ell-\slashed k_2+m_qc)
\right]}
{D_0D_1D_2},
\end{align*}
其中
\begin{align*}
D_0&=\ell^2-m_q^2c^2+\ii0,\\
D_1&=(\ell+k_1)^2-m_q^2c^2+\ii0,\\
D_2&=(\ell-k_2)^2-m_q^2c^2+\ii0.
\end{align*}
第二种排序由 $(a,\mu,k_1)\leftrightarrow(b,\nu,k_2)$ 得到。颜色部分立即给
\begin{equation*}
\Tr(T^aT^b)+\Tr(T^bT^a)
=2T_F\delta^{ab}=\delta^{ab},
\end{equation*}
所以一个颜色单态 Higgs 只能产生 $\delta^{ab}$ 颜色结构。

现在真正做圈积分。三个分母用 Feynman 参数合并：
\begin{equation*}
\frac1{D_0D_1D_2}
=2\int_0^1\dd x\int_0^{1-x}\dd y\,
\frac1{[xD_1+yD_2+(1-x-y)D_0]^3}.
\end{equation*}
令
\begin{equation*}
L=\ell+xk_1-yk_2,
\end{equation*}
则
\begin{equation*}
xD_1+yD_2+(1-x-y)D_0
=L^2-\Delta_q(x,y)+\ii0,
\end{equation*}
其中
\begin{equation*}
\Delta_q(x,y)=m_q^2c^2-xy\,m_h^2c^2.
\end{equation*}
因此阈值结构已经显式出现：当参数区域内 $\Delta_q=0$ 时，对应 $m_h=2m_q$ 的实夸克对分支点。

分子不能只保留质量项。记 $M_q=m_qc$ 以及
\begin{equation*}
a=\ell,\qquad b=\ell+k_1,\qquad c=\ell-k_2.
\end{equation*}
第一张图的 Dirac 迹为
\begin{align*}
N^{\mu\nu}
&=\Tr[(\slashed a+M_q)\gamma^\mu
(\slashed b+M_q)\gamma^\nu(\slashed c+M_q)].
\end{align*}
奇数个 $\gamma$ 的项为零，只剩一个 $M_q$ 与三个 $M_q$ 的项。逐项使用四 $\gamma$ 迹公式可得
\begin{align*}
\frac{N^{\mu\nu}}{4M_q}
={}&b^\mu c^\nu+c^\mu b^\nu-\eta^{\mu\nu}b\!\cdot c
+a^\mu c^\nu-a^\nu c^\mu+\eta^{\mu\nu}a\!\cdot c\\
&+a^\mu b^\nu+a^\nu b^\mu-\eta^{\mu\nu}a\!\cdot b
+M_q^2\eta^{\mu\nu}.
\end{align*}
平移到 $L$ 后，奇数次 $L$ 积分消失，秩二积分按
\begin{equation*}
\int\dd^dL\,L^\mu L^\nu F(L^2)
=\frac{\eta^{\mu\nu}}{d}\int\dd^dL\,L^2F(L^2)
\end{equation*}
约化。两张排序图相加时，非横向的局域二点型部分精确相消；剩余张量必满足
\begin{equation*}
k_1^\mu\mathcal M_{\mu\nu}^{ab}=0,
\qquad
k_2^\nu\mathcal M_{\mu\nu}^{ab}=0,
\end{equation*}
所以可以写成
\begin{equation*}
\mathcal M_{q}^{ab,\mu\nu}
=\delta^{ab}C_q
\left(k_1\!\cdot k_2\,\eta^{\mu\nu}-k_2^\mu k_1^\nu\right).
\end{equation*}
把显式分子投影到这个横向张量后，所有圈动量积分化为一个参数积分：
\begin{equation*}
A_{1/2}(\tau_q)
=4\tau_q\int_0^1\dd x\int_0^{1-x}\dd y\,
\frac{1-4xy}{\tau_q-4xy-\ii0},
\qquad
\tau_q\equiv\frac{4m_q^2}{m_h^2}.
\end{equation*}
这一步与前文 $h\to\gamma\gamma$ 的费米子环具有相同 Lorentz 积分，但这里的电荷因子 $N_cQ_q^2g_{\rm em}^2$ 被 QCD 颜色迹与 $g_s^2$ 取代。令
\begin{equation*}
I(\tau)=\int_0^1\dd x\int_0^{1-x}\dd y\,
\frac1{\tau-4xy-\ii0},
\end{equation*}
先积掉 $y$：
\begin{equation*}
I(\tau)=\int_0^1\frac{\dd x}{4x}
\ln\!\left[\frac{\tau}{\tau-4x(1-x)-\ii0}\right].
\end{equation*}
定义
\begin{equation*}
f(\tau)=2I(\tau)
=\begin{cases}
\arcsin^2(1/\sqrt\tau),&\tau\ge1,\\[4pt]
-\dfrac14\left[
\ln\dfrac{1+\sqrt{1-\tau}}{1-\sqrt{1-\tau}}-\ii\pi
\right]^2,&\tau<1,
\end{cases}
\end{equation*}
再把 $1-4xy=(\tau-4xy)+(1-\tau)$ 分开，得到
\begin{align*}
A_{1/2}(\tau)
&=4\tau\left[\frac12+(1-\tau)I(\tau)\right]\\
&=2\tau\left[1+(1-\tau)f(\tau)\right].
\end{align*}
因此单味系数为
\begin{equation*}
C_q=\frac{\alpha_s}{4\pi v_E}A_{1/2}(\tau_q),
\end{equation*}
对所有夸克味求和即得到正文\eqnrefs{eq:sm-hgg-amplitude-dp94}。这里的 $A_{1/2}$ 不是借用的经验函数，而是三角圈参数积分的闭式。

由振幅进入衰变宽度时，记
\begin{equation*}
C_{gg}=\frac{\alpha_s}{4\pi v_E}\sum_qA_{1/2}(\tau_q),
\qquad
T_{\mu\nu}=k_1\!\cdot k_2\,\eta_{\mu\nu}-k_{2\mu}k_{1\nu}.
\end{equation*}
Ward 恒等式允许用 $-\eta_{\mu\nu}$ 完成物理偏振求和。颜色求和与 Lorentz 收缩分别给
\begin{equation*}
\sum_{a,b}\delta^{ab}\delta^{ab}=N_c^2-1=8,
\qquad
T_{\mu\nu}T^{\mu\nu}=2(k_1\!\cdot k_2)^2
=\frac{m_h^4c^4}{2}.
\end{equation*}
故
\begin{equation*}
\sum|\mathcal M|^2=4|C_{gg}|^2m_h^4c^4.
\end{equation*}
两个末态胶子完全相同，二体相空间必须除以 $2!$：
\begin{equation*}
\Gamma
=\frac1{2!}\frac1{2m_hc^2}\frac1\hbar
\int\dd\Phi_2\sum|\mathcal M|^2,
\qquad
\int\dd\Phi_2=\frac1{8\pi}.
\end{equation*}
再用 $G_F^{(E)}/\sqrt2=1/(2v_E^2)$，便得到正文\eqnrefs{eq:sm-hgg-width-dp94}。

最后才取重顶受控极限。$\tau_t\gg1$ 时
\begin{equation*}
f(\tau_t)=\frac1{\tau_t}+\frac1{3\tau_t^2}+\order(\tau_t^{-3}),
\end{equation*}
所以
\begin{equation*}
A_{1/2}(\tau_t)=\frac43+\frac{14}{45\tau_t}+\order(\tau_t^{-2}).
\end{equation*}
若写局域理论
\begin{equation*}
\Lag_{\rm eff}=C_h\frac{h}{v_E}G^a_{\mu\nu}G^{a\mu\nu},
\end{equation*}
其两胶子树级矩阵元为
\begin{equation*}
\mathcal M_{\mu\nu}^{ab}
=4C_h\frac{\delta^{ab}}{v_E}
\left(k_1\!\cdot k_2\,\eta_{\mu\nu}-k_{2\mu}k_{1\nu}\right).
\end{equation*}
与完整一圈结果的首项逐项匹配：
\begin{equation*}
4C_h\frac1{v_E}
=\frac{\alpha_s}{4\pi v_E}\frac43,
\end{equation*}
从而
\begin{equation*}
\boxed{C_h=\frac{\alpha_s}{12\pi}.}
\end{equation*}
下一项相对修正由 $1/\tau_t=m_h^2/(4m_t^2)$ 控制，所以局域算符不是“把顶夸克删掉”的猜测，而是完整三角圈在低外动量下的系统展开。
\end{derivation}
